\section{Signed refinements of the local comparison}\label{ms:local42-refinements}
The refinements in this appendix apply to an already acquired common-law
resource panel. They preserve the actual count while integrating a Gaussian
increment, and retain the sign of the cubic correction. These are needed
for the proposed comparison with floor $42$; the earlier floor-$47$ formulas
alone do not supply them. Every numerical positivity and coverage statement
below is an acceptance condition. Whenever a normalized binomial mass or
a Stirling expression is used, assume $n\ge N\ge2$ and an interior
integer count $1\le k\le n-1$, so that $\pi_k>0$.

Write $\xi=H/\ell$ for the raw threshold in Bernoulli coordinates. The normalization constants are supplied by
Appendix~\ref{ms:normalization}. The underlying count identities, moment
and tail comparisons are in Section~\ref{ms:geometry-counts} and
Appendix~\ref{ms:fourier:local}.

\subsection{Notation}
Let $Y=B+U$, where $\Pp(B=1)=p$, $\E(U\mid B)=0$,
$v=p(1-p)$, $s=\E U^2$, $\lambda=ns$, and
$\ell=(v+s)^{-1/2}$. The actual threshold has the form
\[
 x=np+H/\ell=k+u,\qquad |u|\le\tfrac12,
 \qquad \delta=k-np=H/\ell-u.
\]
The integer $k$ and the value $\delta$ stay fixed when integrating a function
of a new variable $t$ below. Write $\pi_k=\Pp(\operatorname{Bin}(n,p)=k)$ and
\[
 f(t)=\frac{\Phi((\delta+t)/\sqrt{n(v+s)})-
                   \Phi(\delta/\sqrt{n(v+s)})}{\pi_k}.
\]
The conditional errors satisfy $|U_j|\le A$, $U_0\le a_0$,
$U_1\ge-a_1$, with $a_0,a_1\ge0$ and $d=a_0+a_1$.
All quantities in an application refer to the same iid marginal.

\begin{lemma}[A joint budget on a noise interval]\label{ms:local42:lem:joint}
Suppose the acquired same-law resource inequality is
\[
        \frac d{d_*}+\frac\lambda\Lambda\le1,
        \qquad d_*>0,\quad\Lambda>0.
\]
On a closed noise interval $L\le\lambda\le U\le\Lambda$ one may use
\[
                d\le d_L:=d_*(1-L/\Lambda).
\]
Every conditional support, cubic-moment, tail-distance and characteristic
function bound whose hypothesis is $d\le d_0$ may therefore be applied with
$d_0=d_L$ on that interval.
\end{lemma}
\begin{proof}
Rearrange the displayed resource inequality and use $\lambda\ge L$.
\end{proof}

For example, the whole-count kernel uses $c=1-d_0$. Its moment coefficients
are recomputed as
\[
 \alpha=6m_0\ell_-^2(1-d_0),\qquad
 \beta=\min\{3(4.16)\ell_-(1-p_+-d_0),
                  3(6.16)\ell_-(p_--d_0)\},
\]
where $m_0,\ell_-$ are the previously justified moment and scale lower
bounds. The ordinary cubic correction is recomputed with the same $d_0$.
Likewise, a noise amplitude bound $A$ gives a conditional support span at
most $A+d_0$; the characteristic-function rate is recomputed from that span.
These substitutions are licensed only when the original coefficient
inequalities and positivity gates hold for the acquired moment and scale
bounds.

\subsection{A binomial mass bound on the actual phase cell}
Suppose the finite Stirling argument has supplied positive constants
$E,L_0,R$ on a fixed segment containing every $k/n$ under consideration, with
\[
 \sqrt n\,\pi_k=\frac{\phi(0)}{\sqrt v}
       \exp\{-\delta^2/(2nv)+\zeta\},\qquad
 n|\zeta|\le Z_N(|\delta|),\qquad
 Z_N(t)=\frac{L_0t}{2}+R+\frac{Et^3}{6N},
 \quad n\ge N.
\]
The constants may have been computed on a much wider global phase range.
Let $0<v_-\le v\le v_+$ and let $q_+\ge\sqrt{v_+}$.
\begin{lemma}[Mass localization]\label{ms:local42:lem:mass}
On any smaller closed cell where $|\delta|\le D$,
\[
 \sqrt n\,\pi_k\ge
 \frac{0.39886}{q_+}
 \left[1-\frac{D^2/(2v_-)+Z_N(D)}N\right].
\]
The maximum of this expression and any previously proved mass lower bound
is also a valid mass lower bound. There is no requirement $D\ge1$.
\end{lemma}
\begin{proof}
Lemma~\ref{ms:gaussian-constants} gives $\phi(0)>0.39886$. Set $A_*=D^2/(2v_-)+Z_N(D)\ge0$. The logarithmic formula gives
$\sqrt n\pi_k\ge(0.39886/q_+)e^{-A_*/n}$.
Now $e^{-z}\ge1-z$ and $n\ge N$. If the bracket is negative the displayed
inequality still holds, and the previously positive bound can be retained.
The global Stirling segment, rather than the smaller value $D$, continues
to justify $E,L_0,R$.
\end{proof}

\subsection{A sharper elementary harmonic constant}
\begin{lemma}\label{ms:local42:lem:cos}
Assume the finite rational positivity condition in the proof below.
Then, for every $x\ge0$,
\[
               1-\cos(2\pi x)\le\frac{4553}{1000}x.
\]
\end{lemma}
\begin{proof}
Put $K=4553/1000$. For $x\ge2/K$ the assertion follows from
$1-\cos\le2$. On $[0,2/K]$ take the degree-fourteen alternating upper
Taylor polynomial $Q(x)$ for $1-\cos(2\pi x)$. Its remainder is nonpositive:
the first omitted term is negative, and the remaining term magnitudes
decrease, since $(4\pi/K)^2<16\cdot15$.

Use the directed endpoints
\[
 \pi_-=3.141592653589793<\pi<3.141592653589794=\pi_+.
\]
In each positive coefficient of $Q$ replace $\pi$ by $\pi_+$; in each
negative coefficient replace it by $\pi_-$. This only increases $Q$.
Expand $P(x)=K-Q(x)/x$, a polynomial of degree thirteen,
on each of the $64$ intervals $[2j/(64K),2(j+1)/(64K)]$. If
$P(x)=\sum_{h=0}^{13}p_hx^h$, its Bernstein coefficients on $[a,a+w]$ are
\[
 b_k=\sum_{i=0}^k\frac{\binom ki}{\binom{13}i}
        \sum_{h=i}^{13}p_h\binom hi a^{h-i}w^i,
               \qquad k=0,\ldots,13.
\]
The remaining certificate condition is that all 896 rational coefficients
are positive. Under that condition,
$P>0$ on the full interval. The value at zero follows by
continuity (or directly).

For completeness the stated $\pi$ endpoints are checked using alternating
series for $\arctan(1/5)$ (32 terms) and $\arctan(1/239)$ (12 terms).
Machin's identity was established in Lemma~\ref{ms:gaussian-constants};
the term counts here give the finer directed endpoints required by this
polynomial. The finite recipe uses only rational operations and the Bernstein identity
proved in Section~\ref{ms:certificates}.
\end{proof}

If the original characteristic-function argument gives
$|\E e^{2\pi iW}|\le e^{-rt}$, Lemma~\ref{ms:local42:lem:cos} upgrades its fractional
part estimate from $(5/23)(1-e^{-rt})$ to
\[
                 \frac{1000}{4553}(1-e^{-rt}).
\]
Consequently every occurrence of $5/23$ and $5/46$ in that harmonic gain is
replaced by $1000/4553$ and $500/4553$, respectively. The existing degree-six
Taylor gain and its monotonicity argument remain unchanged.

\subsection{The Gaussian derivative knows the sign of the threshold}
\begin{lemma}[Positive phase, with a fixed actual count]\label{ms:local42:lem:mid}
Assume $u\in[a,b]\subset[0,1/2]$, $H/\ell\le \xi_*$, $|\delta|\le D$,
$v\ge v_->0$, $\lambda\le U$, and the Stirling remainder bound above.
Set
\[
 G=\max\left(0,\max_{z\in[a,b]}(2\xi_*z-z^2)\right),\qquad
 \eta=\frac{U+G}{2v_-}+Z_N(D).
\]
Then $f(u)\ge u(1-\eta/N)$ for every $n\ge N$.
The inner maximum is attained at the projection of $\xi_*$ onto $[a,b]$.
If $H\le0$, one may take $\xi_*=0$, hence $G=0$.
\end{lemma}
\begin{proof}
For a fixed actual law and threshold, differentiate while keeping
$\delta=k-np$ fixed:
\[
 \log f'(t)=-\tfrac12\log(1+s/v)
       +\frac{\delta^2}{2nv}-\frac{(\delta+t)^2}{2n(v+s)}-\zeta.
\]
For $0\le t\le u$, use $\log(1+s/v)\le s/v$ and $v+s\ge v$ to obtain
\[
 \log f'(t)\ge-\frac{\lambda+2\delta t+t^2}{2nv}-\zeta.
\]
The convex function $t\mapsto2\delta t+t^2$ attains its maximum on
$[0,u]$ at an endpoint. Since the actual relation is $\delta=H/\ell-u$,
\[
 2\delta t+t^2\le\max(0,2(H/\ell)u-u^2)\le G.
\]
Thus $\log f'(t)\ge-\eta/n$, and
$f'(t)\ge e^{-\eta/n}\ge1-\eta/N$. Integrate from zero to the actual $u$.
\end{proof}

A signed $H$ interval supplies the bound $D$. If $H\ge0$,
one can use $\xi_*=H_+/\ell_-$ and $D=\max(b,\xi_*-a)$; if $H\le0$,
use $D=H_-/\ell_-+b$. The new lower derivative constant is maximized with
the older valid constant. Every resulting bound is checked positive.

The normal variance and cubic-moment repair is localized separately.
On a phase cell, use its actual upper bound for $\delta_+$, rather than
$|\delta|$. If the resulting coefficient is $\kappa$, replace it by
$\max(0,\kappa)$ before using $\kappa/n\le\kappa/N$.
That sign condition is essential for an all-index conclusion.

\subsection{Reusing a curvature bound}
For fixed probability, threshold and phase intervals, the curvature
function in Section~\ref{ms:curvature} is nonincreasing in $c$.
Lemma~\ref{ms:local42:lem:joint} makes $c=1-d_L$ increase as the lower
noise endpoint $L$ increases. A curvature upper bound for the first band
therefore remains valid on later bands. The other $c$-dependent moment
and support bounds are evaluated using the value for each band.

\subsection{Signed Gaussian corrections on both phase halves}
Let $n\ge N$, $v=p(1-p)>0$, $s\ge0$, $\lambda=ns$, and let the
actual nearest count and phase satisfy
\[
 \delta=k-np=\xi-u,\qquad
 f(t)=\frac{\Phi((\delta+t)/\sqrt{n(v+s)})-
                   \Phi(\delta/\sqrt{n(v+s)})}{\pi_k}.
\]
The actual $k$ and $\delta$ are held fixed when $t$ varies. Assume the
proved finite Stirling formula on this count range is
\[
 \sqrt n\,\pi_k=\frac{\phi(0)}{\sqrt v}
       \exp\{-\delta^2/(2nv)+\zeta\},\qquad
 n|\zeta|\le Z_N(D),\quad |\delta|\le D.
\]
Write $v\ge v_->0$ and $Z=Z_N(D)\ge0$.

\begin{lemma}[Negative phase]
Suppose $u=-r$, $0\le a\le r\le b$, and $\xi\le \xi_*$. Define
\[
 t_*=\max(0,\min(b,\xi_*+b)),\qquad G_-=2(\xi_*+b)t_*-t_*^2,\qquad
 \eta_-=G_-/(2v_-)+Z.
\]
If $D^2\le Nv_-$ and $\eta_-<N$, then
\[
 f'(-t)\le\frac1{1-\eta_-/N}\quad(0\le t\le r),\qquad
 u-f(u)\le\frac{b\eta_-}{n(1-\eta_-/N)}.
\]
\end{lemma}
\begin{proof}
Direct differentiation and the Stirling formula give
\[
 \log f'(-t)=-\tfrac12\log(1+s/v)
 +\frac{\delta^2s}{2nv(v+s)}
 +\frac{2\delta t-t^2}{2n(v+s)}-\zeta.
\]
Since $\log(1+x)\ge x/(1+x)$ for $x\ge0$, the first two terms are at most
\[
 \frac{s}{2(v+s)}\left(\frac{\delta^2}{nv}-1\right)\le0.
\]
This cancellation removes the variance term from the upper derivative
bound. Because $\delta=\xi+r$, $r\le b$ and $0\le t\le r\le b$,
\[
 2\delta t-t^2\le 2(\xi_*+b)t-t^2\le G_-.
\]
The last maximum is the concave quadratic maximum on $[0,b]$, achieved
at the displayed projection $t_*$. It is nonnegative because $t=0$ is
available. Therefore $\log f'(-t)\le\eta_-/n$.
Using $e^x\le(1-x)^{-1}$ for $0\le x<1$ gives
\[
 f'(-t)-1\le\frac{\eta_-/n}{1-\eta_-/n}
             \le\frac{\eta_-}{n(1-\eta_-/N)}.
\]
Integrate $u-f(u)=\int_0^r(f'(-t)-1)\,dt$.
\end{proof}

Equivalently, $G_-=b^2+2\xi_*b$ for $\xi_*\ge0$, and
$G_-=\max(0,\xi_*+b)^2$ for $\xi_*<0$. This sharper vertex is used in the
whole-count comparison on the left probability panels at floor $42$.

\begin{lemma}[Positive phase]
Suppose $u\in[a,b]\subset[0,1/2]$, $\xi\le \xi_*$, and define
\[
 G_+=\max\left(0,\max_{z\in[a,b]}(2\xi_*z-z^2)\right).
\]
Then
\[
 u-f(u)\le\frac b n
       \left(\frac\lambda{2v_-}+\frac{G_+}{2v_-}+Z\right).
\]
The inner maximum is at the projection of $\xi_*$ onto $[a,b]$.
\end{lemma}
\begin{proof}
The exact derivative identity and $\log(1+s/v)\le s/v$ give
\[
 \log f'(t)\ge-\frac{\lambda+2\delta t+t^2}{2nv}-\zeta.
\]
For fixed actual $u$, the function $2\delta t+t^2$ is convex on $[0,u]$.
Its maximum is at an endpoint, and $\delta=\xi-u$ gives
\[
 2\delta t+t^2\le\max(0,2\xi u-u^2)\le G_+.
\]
Use $e^{-x}\ge1-x$ and integrate from zero to $u$.
\end{proof}

\subsubsection{Whole-count insertion and interval coverage}
The whole-count identity is split as
$Q/\pi_k-f(u)=(Q/\pi_k-u)+(u-f(u))$.
Thus one needs an upper bound on the signed correction $u-f(u)$;
an absolute-value estimate is unnecessary. After division by $\lambda$,
the new coefficients replacing the old Gaussian correction are
\[
 \begin{array}{ll}
 u\ge0:&\displaystyle
 \frac b{2v_-}+\frac{b\{G_+/(2v_-)+Z\}}\lambda,\\[6pt]
 u\le0:&\displaystyle
 \frac{b\eta_-}{(1-\eta_-/N)\lambda}.
 \end{array}
\]
On each fixed phase cell these have the form $A+B/\lambda$ with
$A,B\ge0$. They therefore preserve the original convex endpoint check on
each closed noise interval. Every phase cell is confined to one half-line;
zero is a shared closed endpoint. The count displacement bound $D$ comes
from the same actual $\xi-u$ cell, and the Stirling constants continue to
come from the already proved global count segment.

The other all-index steps are unchanged: the event bracket is checked
negative before substituting a binomial mass lower bound, and the entire
inverse-$n$ correction is checked positive before replacing $n$ by $N$.
The inverse correction is $I(\lambda)=K+A+B/\lambda$, hence decreases with $\lambda$; positivity at both endpoints (in particular the upper endpoint) licenses the all-index step throughout the interval. The final score is convex in $\lambda$: its event fraction is a positive constant divided by an affine positive function, its $B/\lambda$ coefficient is nonnegative, and its remaining terms are affine or constant. These are independent gates in the exact kernel. The formulas apply to
both signs of $\xi$ without reflecting the law or changing the count during
integration. They also supply an improved upper derivative for the
negative-phase nearest-count inequalities.
\subsubsection{A complete phase cover below the whole-count threshold}
For $\lambda\in[0.05,0.08]$, the shifted-quartic nearest-count comparisons
must be checked on every closed phase interval of length $0.001$ in
$[0,1/2]$, for both phase signs. The shifted-quartic polynomial and
nearest-count identity are defined at phase zero, so the two closed covers
include $u=0$ directly. This permits the whole-count comparison to
start at $\lambda=0.08$. The four nearest-count bands use phase starts
$0,0,0.19,0.27$ at noise bands $[0.05,0.06]$, $[0.06,0.08]$,
$[0.08,0.10]$, and $[0.10,0.12]$, respectively.

\subsection{Conditional centering and the exact cubic bound}
Let $B$ be Bernoulli with parameter $0<p<1/2$, put $q=1-p$ and $v=pq$,
and write $Y=B+U$, with $\E(U\mid B)=0$.
Let $U_0,U_1$ denote the conditional laws of $U$, and assume
\[
 U_0\le a_0<p,\qquad U_1\ge-a_1>-(1-p),\qquad
 a_0,a_1\ge0,\qquad d=a_0+a_1.
\]
Put $s_j=\E U_j^2$, $s=qs_0+ps_1$, $\lambda=ns$, and
\[
 X=\frac{Y-p}{\sqrt{v+s}},\quad
 \rho=\E|X|^3,\quad
 \rho_B=\frac{1-2v}{\sqrt v}.
\]
Finite third moments are assumed. The support conditions fix the sign of
$Y-p$ in each class, which is necessary for the following expansion.

\begin{lemma}[No positivity assumption on the cubic bracket]\label{ms:local42:cubic}
With the preceding notation,
\begin{equation}\label{ms:local42:bracket}
 (v+s)^{3/2}(\rho-\rho_B)\ge sQ,\qquad
 Q=6p-3p^2-\frac32-d-\frac{3(1-2v)s}{8v}.
\end{equation}
\end{lemma}
\begin{proof}
Expanding within each sign class and using conditional centering gives
\[
 \E|Y-p|^3=v(1-2v)+3pq(s_0+s_1)
                         +p\E U_1^3-q\E U_0^3.
\]
The one-sided support bounds imply
$\E U_1^3\ge-a_1s_1$ and $\E U_0^3\le a_0s_0$.
Moreover $3pq(s_0+s_1)\ge3p(qs_0+ps_1)=3ps$, because $q\ge p$.
Consequently
\[
                         \E|Y-p|^3\ge v(1-2v)+(3p-d)s.
\]
For $z\ge0$, Taylor's formula gives
$(1+z)^{3/2}\le1+3z/2+3z^2/8$: the third derivative is negative.
Subtract
\[
 \rho_B(v+s)^{3/2}
 =v(1-2v)(1+s/v)^{3/2}
 \le v(1-2v)+\tfrac32(1-2v)s+\frac{3(1-2v)s^2}{8v}
\]
from the preceding moment lower bound and use
$3p-\tfrac32(1-2v)=6p-3p^2-\tfrac32$.
\end{proof}

\subsection{The coefficient must be bounded in the correct direction}
Suppose $n\ge N$, $p\in[p_-,p_+]\subset(0,1/2)$,
$0\le d\le d_0$, and $0\le\lambda\le\Lambda$.
Choose $0<v_*\le p_-(1-p_-)$ and put
\[
 V^*=p_+(1-p_+)+\Lambda/N,\qquad
 Q_*=6p_--3p_-^2-\frac32-d_0
                     -\frac{3(1-2v_*)\Lambda}{8v_*N}.
\]
Then $Q\ge Q_*$: the first term increases with $p$, and
$(1-2v)/v$ decreases with $v$. In particular, $Q_*$ may be negative.
Already with $d=s=0$ its sign changes at $p=1-1/\sqrt2$.

Let $C$ be the score multiplying the moment correction, and assume
$0<c_*\le C\le c^*$. Choose positive multipliers satisfying
\[
        0<m_-\le\frac{c_*}{(V^*)^{3/2}},\qquad
        m_+\ge\frac{c^*}{v_*^{3/2}},
\]
and define the signed gain
\begin{equation}\label{ms:local42:gain}
             G_*:=\min\{m_-Q_*,m_+Q_*\}.
\end{equation}

\begin{proposition}[Signed gain and nonnegative repair]\label{ms:local42:repair}
Let $\Phi$ denote the standard normal distribution function. For an actual
count displacement $\delta$, set
\[
 \mathcal R=\sqrt n\left\{
 \Phi\!\left(\frac{\delta}{\sqrt{nv}}\right)
 -\Phi\!\left(\frac{\delta}{\sqrt{n(v+s)}}\right)\right\}
                         +C(\rho_B-\rho).
\]
If $D_+\ge\max(\delta,0)$ and $P_*\ge 0.4/(2v_*^{3/2})$, then
\begin{equation}\label{ms:local42:repair_bound}
 \mathcal R\le(P_*D_+-G_*)\frac{\lambda}{n}
 \le \kappa\frac{\lambda}{N},\qquad
 \kappa=\max\{0,P_*D_+-G_*\}.
\end{equation}
\end{proposition}
\begin{proof}
Write $A=C/(v+s)^{3/2}$. The hypotheses give
$0<m_-\le A\le m_+$. By Lemma~\ref{ms:local42:cubic},
$C(\rho-\rho_B)\ge sAQ\ge sAQ_*$.
If $Q_*\ge0$, $AQ_*\ge m_-Q_*$; if $Q_*<0$,
$AQ_*\ge m_+Q_*$. This proves $C(\rho-\rho_B)\ge sG_*$
in both cases. One must use the \emph{upper} multiplier in the negative
case.

The normal part is nonpositive when $\delta\le0$.
For $\delta>0$, differentiation with respect to an intermediate variance
$t\in[0,s]$ gives
\[
 \frac{d}{dt}\left[-\sqrt n\,
 \Phi\!\left(\frac{\delta}{\sqrt{n(v+t)}}\right)\right]
 =\frac{\delta}{2(v+t)^{3/2}}
   \phi\!\left(\frac{\delta}{\sqrt{n(v+t)}}\right)
 \le\frac{0.4\delta}{2v_*^{3/2}}\le P_*D_+.
\]
Integrate over $[0,s]$, and combine with the moment bound.
Finally $s=\lambda/n$ and $\kappa\ge0$, so
$(P_*D_+-G_*)\lambda/n\le\kappa\lambda/n\le\kappa\lambda/N$.
The clamp at zero is needed before replacing $n$ by its lower bound.
\end{proof}

\subsection{Exact checks and the finite application}
The left probability panels at floor $42$ use
\[
                    N=42,\qquad c_*=0.4,\qquad c^*=0.4226.
\]
These are hypotheses supplied by the surrounding finite-index proof.
The required rational multiplier inequalities are
\[
 (m_-)^2(V^*)^3<c_*^2,\qquad
 (m_+)^2v_*^3>(c^*)^2,
\]
\emph{together with positivity of both multipliers}. Squaring without
those positivity checks would accept a false negative coefficient.
The normal multiplier is checked similarly using
$4P_*^2v_*^3>0.4^2$ and $P_*>0$.

On each probability and noise interval, evaluate the minimum in
\eqref{ms:local42:gain} and bound the nonnegative coefficient in
\eqref{ms:local42:repair_bound} using the same moment and support data.
Lemma~\ref{ms:local42:cubic} and Proposition~\ref{ms:local42:repair}
then give the probabilistic correction.

\subsection{The finite comparison cover at floor $42$}\label{ms:local42-cover}
The proposed acquisition first excludes the closed probability ranges
$[0,0.23]$, $[0.23,0.29]$, $[0.474,0.48]$ and $[0.48,0.63]$, together with
$[0.63,1]$. Different bounds may establish adjacent exclusions, so their
common endpoints must be retained. The remaining grids are the $45$ left
panels in $[0.29,0.38]$ and $47$ central panels in $[0.38,0.474]$, each of width
$0.002$. These are the rows defined in Section~\ref{ms:resources}.
The impossible endpoints $p=0,1$ are excluded by centering and unit variance.
Each exclusion and each resulting moment interval is a certificate obligation.

On a surviving row, intersect every noise interval with its acquired
ceiling $\Lambda$. The following table specifies the original comparison
cover. A whole-count radius $r$ means $|u|\le r$; the nearest-count
comparison covers the complementary closed phase intervals up to $|u|=1/2$. Both signs of $H$ are checked
separately. At $\lambda=0$, the actual law is Bernoulli and the separately
cited binomial theorem applies.
\begin{center}\small
\begin{tabular}{@{}lP{55mm}P{55mm}@{}}\toprule
Noise interval & Central panels & Left panels\\\midrule
$(0,0.05]$ & Small-noise bound, all phases & Small-noise bound, all phases\\
$[0.05,0.06]$ & Whole count: $|u|\le0.09$; nearest count outside & Nearest count, all phases\\
$[0.06,0.08]$ & Whole count: $|u|\le0.14$; nearest count outside & Nearest count, all phases\\
$[0.08,0.10]$ & Whole count: $|u|\le0.19$; nearest count outside & Whole count: $|u|\le0.19$; nearest count outside\\
$[0.10,0.12]$ & Whole count: $|u|\le0.29$; nearest count outside & Whole count: $|u|\le0.27$; nearest count outside\\
$[0.12,\Lambda]$ & Whole count, all phases & Whole count, all phases\\\bottomrule
\end{tabular}
\end{center}
All seams are closed and shared. The comparison's amplitude caps must also
be acquired: on the four bands ending at $0.08,0.10,0.12,\Lambda$, the
central-panel caps are $0.375,0.325,0.275,0.25$ and the left-panel caps are
$0.40,0.35,0.30,0.275$, respectively, intersected with the row's global cap.
A narrower cap is an output of confinement before it can be used here.
The joint gap cap is recomputed from the same weighted resource in every
band. The additional exclusions of Section~\ref{ms:resources} reduce this
cover once their required moment and resource bounds have been established.
